\subsection{等价关系的定义}

让 \(R\{x, y\}\) 是一个关系，其中x和y是不同的字母。关系R被称为（关于字母x和y）对称的，如果

\[
\mathrm{R} \{x, y \} \Longrightarrow \mathrm{R} \{y, x \}.
\]

如果是这样，通过将x和y替换为两个字母 \(x'\) 并且 \(y'\) ，它们彼此不同且与R中出现的所有字母都不同，然后依次将y和x替换为 \(x'\) 并且 \(y'\) ，我们可以看到

\[
\mathrm{R} \left\{y, x \right\} \Longrightarrow \mathrm{R} \left\{x, y \right\}.
\]

因此， \(\mathbb{R}\{x, y\}\) 并且 \(\mathbb{R}\{y, x\}\) 是等价的。设 \(z\) 为一个未出现在 \(\mathbb{R}\) 。关系 \(\mathbb{R}\) 被称为可传递的（关于字母） \(x, y\) ）如果我们有

\[
(\mathrm{R} \{x, y \} \text {   且   } \mathrm{R} \{y, z \}) \Longrightarrow \mathrm{R} \{x, z \}.
\]

例子。关系 x = y 是对称且传递的。关系 \(X \subset Y\) 是传递的但不对称。关系 \(X \cap Y = \varnothing\) 是对称的但不传递的。

如果 \(\mathbf{R}\{x, y\}\) 既对称又传递，就称它是（关于字母 \(x\) 和 \(y\) 的）等价关系。在这种情况下，记号 \(x \equiv y \pmod{\mathbf{R}}\) 有时用作 \(\mathbf{R}\{x, y\}\) 的同义记号；读作“ \(x\) 等价于 \(y\) （模 \(\mathbf{R}\) ）”。如果 \(\mathbf{R}\) 是等价关系，则有 \(\mathbf{R}\{x, y\} \Longrightarrow (\mathbf{R}\{x, x\} \text{ 且 } \mathbf{R}\{y, y\})\) ，因为 \(\mathbf{R}\{x, y\}\) 蕴含 \(\mathbf{R}\{y, x\}\) ，而由定义， \((\mathbf{R}\{x, y\} \text{ 且 } \mathbf{R}\{y, x\})\) 蕴含 \((\mathbf{R}\{x, x\} \text{ 且 } \mathbf{R}\{y, y\})\) 。

让 \(\mathbf{R}\{x, y\}\) 是一个关系；它被称为在 \(\mathbf{E}\) （就字母而言） \(x, y\) ) 如果该关系 \(\mathbf{R}\{x, x\}\) 等价于 \(x \in \mathbf{E}\) 。如果关于 \(\mathbf{E}\) 不存在可能的歧义，则简言之（滥用语言地）称 \(\mathbf{R}\) 是自反的。

¶ E上的等价关系被定义为在E上自反的等价关系。如果R\{x, y\} 是 E 上的一个等价关系，我们有 R\{x, y\} \(\Longrightarrow\) ((x, y) ∈ E × E)，因此 R 有一个图（关于字母 x, y）。反之，假设等价关系 R\{x, y\} 有一个图 G。注意关系 R\{x, x\} 等价于关系 (∃y)R\{x, y\}；因为前者蕴含后者（第一章，§4，第2节，方案S5），反之亦然，由于R\{x, y\} 蕴含R\{x, x\}，(∃y)R\{x, y\} 意味着存在某个y使得R成立\{x, x\} 因此也有R\{x, x\}。因此R\{x, x\} 等价于 \(x \in pr_{1}G\) ，且因此R是 \(pr_{1}G\) .

集合 E 上的一个等价关系是一个对应，其源与目标都等于 E，且其图形 F 使得关系 \((x, y) \in \mathbb{F}\) 是 E 上的一个等价关系。

\minorheading{示例}

\begin{enumerate}
\def\labelenumi{(\arabic{enumi})}
\item
  关系 \(x = y\) 是一个没有图形的等价关系，因为如果它有图形，则该图形的第一投影将是所有对象的集合。
\item
  关系“x = y 且 \(x \in E\) '' 是 E 上的一个等价关系，其图形是 \(E \times E\) .
\item
  关系“存在一个从 X 到 Y 的双射”是一个等价关系，但它没有图形（参见第三章，第 3 节）。
\item
  关系“ \(x \in E\) 并且 \(y \in E\) '' 是 \(E\) 上的一个等价关系，其图形是 \(E \times E\) .
\item
  设 \(\mathbf{A} \subset \mathbf{E}\) ；则关系
\end{enumerate}

\[
(x \in \mathrm{E} - \mathrm{A} \text { 且 } y = x) \text { 或 } (x \in \mathrm{A} \text { 且 } y \in \mathrm{A})
\]

是 E 上的一个等价关系。

\begin{enumerate}
\def\labelenumi{(\arabic{enumi})}
\setcounter{enumi}{5}
\tightlist
\item
  * 关系“ \(x \in \mathbf{Z}\) 并且 \(y \in \mathbf{Z}\) 并且 \(x - y\) 能被4整除”是 \(\mathbf{Z}_{*}\) .
\end{enumerate}

命题1. 一个对应关系 \(\Gamma\) X 与 X 之间的等价关系当且仅当它满足以下条件：(a) X 是 \(\Gamma\) 的定义域；(b) \(\Gamma = \overline{\Gamma}^1\) ；(c) \(\Gamma \circ \Gamma = \Gamma\) .

设 \(\Gamma\) 是 X 与 X 之间的对应关系，并设 G 为其图像。如果 \(\Gamma\) 是 X 上的一个等价关系，则 \((x, x) \in G\) 对于所有 \(x \in X\) ；因此 X 是 \(\Gamma\) 的定义域。关系 \((x, y) \in G\) 等价于 \((y, x) \in G\) ，因此等价于 \((x, y) \in \overline{G}\) ，因此 \(G = \overline{G}\) ，从而 \(\Gamma = \overline{\Gamma}\) 。关系 \((x, y) \in G\) 和 \((y, z) \in G\) 蕴含 \((x, z) \in G\) ，因此 \(G \circ G \subset G\) ；反之， \((x, y) \in G\) 蕴含 \((x, x) \in G\) ，因此 \((x, y) \in G \circ G\) ，因此 \(G \subset G \circ G\) ；因此 \(G = G \circ G\) ，从而 \(\Gamma = \Gamma \circ \Gamma\) .

反之，假设条件（a）、（b）和（c）均满足。该关系 \((x, y) \in G\) 是对称的（由(b)）且传递的（由(c)）；因此它是一个等价关系，并且由(a)可知它是上的一个等价关系。 \(X\) .

